\section*{Capítulo \ref{ch:b2:oceanic-islands} --- Ilhas oceânicas}

\begin{solution}{exo:b2:oceanic-islands:1}
(a) Fase pré-escudo, como o Loihi hoje. (b) \omterm{def:b2:oceanic-islands:stages}{Fase de escudo}. (c) Fase pós-escudo. (d) Fase
rejuvenescida. (e) Nenhuma fase eruptiva mais: o vulcão afundou sob seu recife, que agora é um \omterm{def:b2:oceanic-islands:reefs}{atol}.
\end{solution}

\begin{solution}{exo:b2:oceanic-islands:2}
(a) $2500 + 350\sqrt9 = \qty{3550}{m}$. (b) $350\times(\sqrt{25} - \sqrt9) = 350\times2 =
\qty{700}{m}$. (c) $350/(2\times5) = \qty{35}{m}$ por milhão de anos.
\end{solution}

\begin{solution}{exo:b2:oceanic-islands:3}
(a) $5000/17 = \qty{294}{kyr}$ em torno das ilhas principais; $30\,000/70 = \qty{429}{kyr}$ ao longo da
dorsal. (b) Há deslizamentos que escapam: um vulcão mais jovem pode soterrar a cicatriz e os detritos de um mais antigo,
e depósitos antigos são cobertos por sedimento. O número verdadeiro é maior e o intervalo verdadeiro,
mais curto.
\end{solution}

\begin{solution}{exo:b2:oceanic-islands:4}
(a) $4158\times0.3048 = \qty{1267}{m}$; $4610\times0.3048 = \qty{1405}{m}$; $150\times0.3048 =
\qty{46}{m}$. (b) O primeiro recife não cresceu a mais de \qty{46}{m} de profundidade e sua base está hoje a
\qty{1267}{m}: o vulcão afundou pelo menos $1267 - 46 = \qty{1221}{m}$ (mais, se o nível do mar
estava mais alto então). (c) O afundamento é no máximo de cerca de \qty{1267}{m}, ao longo de pelo menos \qty{33.9}{Myr}:
no máximo $1267/33.9 = \qty{37}{m}$ por milhão de anos, \qty{0.037}{mm} por ano.
\end{solution}

\begin{solution}{exo:b2:oceanic-islands:5}
(a) $\dot S = 175/\sqrt t$: \qty{87.5}{m.Myr^{-1}} (\qty{0.0875}{mm.yr^{-1}}) a \qty{4}{Myr},
\qty{35}{m.Myr^{-1}} (\qty{0.035}{mm.yr^{-1}}) a \qty{25}{Myr}, \qty{25}{m.Myr^{-1}}
(\qty{0.025}{mm.yr^{-1}}) a \qty{49}{Myr}. (b) $4.2/0.0875 = 48$ vezes. (c) $2.6/0.0875 = 30$
vezes: a fase de carga faz uma ilha afundar muito mais depressa que o resfriamento de uma placa jovem.
\end{solution}

\begin{solution}{exo:b2:oceanic-islands:6}
(a) $\dot S = (3200/62.8)\,e^{-100/62.8} = 50.96\times0.2035 = \qty{10.37}{m}$ por milhão de anos,
\qty{0.0104}{m} por mil anos: os \qty{0.01}{m} por mil anos medidos. (b) $e^{-t/62.8}
= 10/50.96$, $t = 62.8\ln5.096 = \qty{102.3}{Myr}$. (c) Lei da raiz: $2500 + 350\times5 =
\qty{4250}{m}$; lei exponencial: $6400 - 3200\times0.6716 = \qty{4251}{m}$; elas diferem em
\qty{1}{m}.
\end{solution}

\begin{solution}{exo:b2:oceanic-islands:7}
(a) $\qty{8.6}{cm.yr^{-1}} = \qty{86}{km.Myr^{-1}}$; $86\times33 = \qty{2838}{km}$. (b) Sim:
\qty{27.7}{Myr} é menos que os \qty{33}{Myr} depois dos quais os recifes se afogam, de modo que seu recife deveria
ainda estar vivo. (c) $55\times1.609 = \qty{88.5}{km}$, ou $88.5/86 = \qty{1.0}{Myr}$ a mais: cerca de
\qty{28.7}{Myr}. É um pouco menos que \qty{33}{Myr}, mas os \qty{33}{Myr} são eles próprios um
valor aproximado para a cadeia inteira, e o crescimento depende também da temperatura da água em cada
\omterm{def:b2:oceanic-islands:reefs}{atol}: Kure, o último \omterm{def:b2:oceanic-islands:reefs}{atol} vivo, está de fato perto do \omterm{def:b2:oceanic-islands:darwin-point}{ponto de Darwin}.
\end{solution}

\begin{solution}{exo:b2:oceanic-islands:8}
(a) Caixa: $30\times17\times2 = \qty{1020}{km^3}$; pirâmide: um terço, \qty{340}{km^3}. (b) Um único bloco
é um décimo (caixa) ou um trigésimo (pirâmide) do vulcão inteiro do Kilauea. (c) $23\,000\times
0.1 = \qty{2300}{km^3}$, um quarto do Kilauea; uma estimativa que se sustenta ou cai com a espessura
suposta.
\end{solution}

\begin{solution}{exo:b2:oceanic-islands:9}
(a) $1000/2.6 = \qty{385}{kyr}$. (b) Sim: \qty{2.6}{m} por mil anos são
\qty{2.6}{mm.yr^{-1}}, menos que \qty{4.2}{mm.yr^{-1}}, uma margem de \qty{1.6}{mm.yr^{-1}}.
(c) Subida relativa $12 + 2.6 = \qty{14.6}{mm.yr^{-1}}$; o recife fica para trás $14.6 - 4.2 =
\qty{10.4}{mm.yr^{-1}}$ durante $16.5 - 8.2 = \qty{8.3}{kyr}$: seu topo termina a $10.4\times8.3 =
\qty{86}{m}$. (d) $(12 - 4.2)\times8.3 = \qty{65}{m}$. Nos dois casos ele fica acima do limite de crescimento
mais profundo, mas numa água muito mais escura: se ele recupera o atraso depende de como seu crescimento se desacelera com
a profundidade.
\end{solution}

\begin{solution}{exo:b2:oceanic-islands:10}
(a) $S(t_1) = C(\sqrt{t_1 + \Delta} - \sqrt{t_1})$ tem a derivada $\frac C2(1/\sqrt{t_1 +
\Delta} - 1/\sqrt{t_1}) < 0$: ela diminui com $t_1$. (b) $350\sqrt{20} = \qty{1565}{m}$; $350\times
(\sqrt{70} - \sqrt{50}) = 350\times1.2955 = \qty{453}{m}$. (c) Não: $\sqrt t$ é finita em $0$, de modo que
o afundamento em qualquer intervalo é finito, \qty{1565}{m} em (b). Só a taxa no instante $t = 0$
é infinita, onde a lei já não pretende descrever o fundo em detalhe.
\end{solution}

\begin{solution}{exo:b2:oceanic-islands:11}
(a) $30\times0.1 = \qty{3}{m}$. (b) Cerca de \qty{120}{m}, se seu topo estava ao nível do mar no
nível alto: o \omterm{def:b2:oceanic-islands:reefs}{atol} se torna uma ilha calcária, a chuva dissolve a aragonita dos corais e a
rocha se recristaliza em calcita, com cavernas. (c) $120/3 = 40$ ciclos, \qty{4}{Myr}. (d) As zonas
lixiviadas são esses momentos de exposição, selados sob o recife que voltou a crescer quando o mar voltou.
Ao longo de muitos ciclos, o mar retorna a níveis altos semelhantes, de modo que $\Delta\eta$ é pequeno, o topo do recife
volta a crescer até o nível do mar a cada nível alto, e a espessura da capa é fixada pelo afundamento, como diz a
proposição.
\end{solution}

\begin{solution}{exo:b2:oceanic-islands:12}
(a) $3200\times(e^{-50/62.8} - e^{-120/62.8}) = 3200\times(0.4510 - 0.1480) = \qty{970}{m}$.
(b) $150 + 970 = \qty{1120}{m}$. (c) $700\times1.8288 = \qty{1280}{m}$, \qty{160}{m} mais fundo que a
estimativa. (d) O \omterm{def:b2:oceanic-islands:seamount}{guyot} pode ter se afogado antes, numa placa mais jovem, e afundado mais; um vulcão
vizinho pode tê-lo carregado; seu topo pode ter sido cortado mais baixo pelas ondas, ou estar coberto por sedimento; a
placa sob ele pode ter uma idade diferente; e o nível do mar da época do afogamento não era
o de hoje.
\end{solution}

\begin{solution}{pb:b2:oceanic-islands:1}
\textbf{1.} $2500 + 350\sqrt{10} = 2500 + 350\times3.1623 = \qty{3607}{m}$.
\textbf{2.} $3607 + 1000 = \qty{4607}{m}$.
\textbf{3.} $V = \pi r^2h/3 = \pi\times50^2\times4.607/3 = \qty{12061}{km^3}$.
\textbf{4.} $0.85\times12\,061 = \qty{10252}{km^3}$ em \qty{1}{Myr}: \qty{0.0103}{km^3} por ano,
dez vezes menos que o Kilauea hoje.
\textbf{5.} $95\,600/12\,061 = 7.9$: cerca de oito.
\textbf{6.} Lava toleítica, da \omterm{def:b2:oceanic-islands:stages}{fase de escudo}; as lavas alcalinas da fase pós-escudo o
recobrem.
\textbf{7.} $\qty{1}{km}/\qty{1}{Myr} = \qty{1}{m}$ por mil anos, 2.6 vezes mais lento que ao largo
da ilha do Havaí.
\textbf{8.} Sim: $4.2 - 1.0 = \qty{3.2}{mm.yr^{-1}}$ de folga.
\textbf{9.} $12 + 1 - 4.2 = \qty{8.8}{mm}$ por ano.
\textbf{10.} $150/8.8 = \qty{17}{kyr}$. Um terraço recifal afogado no flanco, um degrau do tipo
encontrado ao largo da ilha do Havaí.
\textbf{11.} $350\times(\sqrt{22} - \sqrt{10}) = 350\times(4.6904 - 3.1623) = \qty{535}{m}$.
\textbf{12.} $1000 - 535 = \qty{465}{m}$, a $465/12 = \qty{38.8}{m}$ por milhão de anos.
\textbf{13.} $350/(2\sqrt{10}) = \qty{55.3}{m.Myr^{-1}}$ no início; $350/(2\sqrt{22}) =
\qty{37.3}{m.Myr^{-1}}$ no fim.
\textbf{14.} $2500 + 350\times4.6904 = \qty{4142}{m}$.
\textbf{15.} As grandes ilhas havaianas se elevam a até cerca de quatro vezes mais alto, de modo que há muito mais
rocha acima do nível do mar a remover; sua \omterm{def:b2:oceanic-islands:stages}{fase de escudo} também durou mais.
\textbf{16.} As ondas, que recuam a costa e aplainam uma plataforma de abrasão; os deslizamentos e os \omterm{def:b2:oceanic-islands:collapse}{colapsos
de flanco}.
\textbf{17.} Um \omterm{def:b2:oceanic-islands:reefs}{recife em barreira}, com uma laguna entre o recife e o que resta da costa vulcânica.
\textbf{18.} $350\times(\sqrt{70} - \sqrt{22}) = 350\times(8.3666 - 4.6904) = \qty{1287}{m}$.
\textbf{19.} $1287/48 = \qty{26.8}{m}$ por milhão de anos, \qty{0.027}{mm} por ano: $4.2/0.0268 =
157$ vezes menos do que o recife poderia crescer.
\textbf{20.} \qty{20}{m} mais espessa que o calcário sobre o basalto sob a ilha Parry
(\qty{1267}{m}), e \qty{118}{m} mais fina que sob Elugelab (\qty{1405}{m}): da mesma ordem.
\textbf{21.} $2500 + 350\times8.3666 = \qty{5428}{m}$.
\textbf{22.} $3200\times(e^{-70/62.8} - e^{-120/62.8}) = 3200\times(0.3280 - 0.1480) =
\qty{576}{m}$.
\textbf{23.} $150 + 576 = \qty{726}{m}$.
\textbf{24.} A capa recifal do \omterm{def:b2:oceanic-islands:reefs}{atol} Petrel tem \textbf{\qty{1287}{m}} de espessura, cerca de \qty{1.3}{km}
de calcário recifal, como sob Enewetak.
\end{solution}
